\chapter*{נספח א: הקוד שמאחורי הגרף}
\addcontentsline{toc}{chapter}{נספח א: הקוד שמאחורי הגרף}
\markboth{נספח א: הקוד שמאחורי הגרף}{}

כל הגרפים בספר נוצרו בקוד \en{Python}, ולא צוירו ביד, וכל הספר נבנה מחדש בפקודה אחת, \en{bash build.sh}. הקטע שלמטה הוא הפונקציה שמייצרת את איור~\ref{fig:goldilocks}. היא מתרגמת את משוואה~\eqref{eq:goldilocks} ישירות לקוד: לכל פעילות היא בוחרת $\beta_1 = 2\beta_2 t^{*}$, כך שהשיא נופל בדיוק על נקודת המפנה שדווחה \parencite{oii2017press, przybylski2017}. שימו לב: מיקום השיא הוא \textbf{קלט} של המודל ולא תוצאה שלו, והגובה מנורמל. הקוד מדגים את הצורה, אבל אינו מוכיח אותה.

\pythonfile[firstline=108, lastline=138]{python/make_graphs.py}{הפונקציה \en{goldilocks()} מתוך \en{python/make\_graphs.py}}
